% GENERATED FILE. Edit canonical lessons, then run npm run generate:books.
% canonical-source-hash: fnv1a64:a9123159f2a34e9e
% canonical-lessons: IT-C244-proprio, IT-C244-quasi, IT-C244-circa, IT-C244-soltanto, IT-C244-insieme

\chapter{Exactly, Almost, About, Only, Together}
\label{ch:it-exactly-almost-about-only-together}
\hlchaptermodality{244}

\begin{chapteropening}
\textbf{By the end of this chapter:} \emph{I can say exactly, almost, about, only and together.}

Complete the last lesson of chapter 244: I can say exactly, almost, about, only and together.
\end{chapteropening}

\section[proprio]{proprio — exactly, really}

\label{lesson:IT-C244-proprio}



\begin{quote}
Before the new one: say the Italian for an emergency, then the Italian for an accident.
\end{quote}

\subsection*{You'll want to know: proprio}
\textbf{proprio} — \textquotedblleft{}exactly, really\textquotedblright{}.

\begin{cousinweb}
\textbf{proprio} $\leftarrow$ Latin \emph{proprius}, \textquotedblleft{}one's own\textquotedblright{} — a learned borrowing. (EtymDB 2.1)
\end{cousinweb}

\begin{grammarlens}[title={What you've built}]
One more word for this chapter.
\end{grammarlens}

\subsection*{Your turn}
\begin{itemize}
\raggedright
  \item \emph{Say it:} \emph{proprio}
  \item \emph{Say it:} \emph{proprio}, once more
  \item \emph{Say it:} say \emph{proprio}
  \item \emph{Recall:} say the Italian for an emergency, then the Italian for an accident, then say \emph{proprio} again
\end{itemize}

\begin{tcolorbox}[breakable,colback=teal!4,colframe=teal!35!black,title={Before you move on}]
What does proprio mean? (\textquotedblleft{}Exactly, really\textquotedblright{}.) Say it once more.
\end{tcolorbox}
\section[quasi]{quasi — almost}

\label{lesson:IT-C244-quasi}



\begin{quote}
Before the new one: say the Italian for an accident, then the Italian for exactly.
\end{quote}

\subsection*{You'll want to know: quasi}
\textbf{quasi} — \textquotedblleft{}almost\textquotedblright{}.

\begin{cousinweb}
\textbf{quasi} $\leftarrow$ Latin \emph{quasi}, \textquotedblleft{}as if\textquotedblright{}, from \emph{quam}, \textquotedblleft{}as\textquotedblright{}, and \emph{si}, \textquotedblleft{}if\textquotedblright{} — inherited. (EtymDB 2.1)
\end{cousinweb}

\begin{grammarlens}[title={What you've built}]
One more word for this chapter.
\end{grammarlens}

\subsection*{Your turn}
\begin{itemize}
\raggedright
  \item \emph{Say it:} \emph{quasi}
  \item \emph{Say it:} \emph{quasi}, once more
  \item \emph{Say it:} say \emph{quasi}
  \item \emph{Recall:} say the Italian for an accident, then the Italian for exactly, then say \emph{quasi} again
\end{itemize}

\begin{tcolorbox}[breakable,colback=teal!4,colframe=teal!35!black,title={Before you move on}]
What does quasi mean? (\textquotedblleft{}Almost\textquotedblright{}.) Say it once more.
\end{tcolorbox}
\section[circa]{circa — about, approximately}

\label{lesson:IT-C244-circa}



\begin{quote}
Before the new one: say the Italian for exactly, then the Italian for almost.
\end{quote}

\subsection*{You'll want to know: circa}
\textbf{circa} — \textquotedblleft{}about, approximately\textquotedblright{}.

\begin{cousinweb}
\textbf{circa} $\leftarrow$ Latin \emph{circa}, \textquotedblleft{}around\textquotedblright{} — a learned borrowing. (EtymDB 2.1)
\end{cousinweb}

\begin{grammarlens}[title={What you've built}]
One more word for this chapter.
\end{grammarlens}

\subsection*{Your turn}
\begin{itemize}
\raggedright
  \item \emph{Say it:} \emph{circa}
  \item \emph{Say it:} \emph{circa}, once more
  \item \emph{Say it:} say \emph{circa}
  \item \emph{Recall:} say the Italian for exactly, then the Italian for almost, then say \emph{circa} again
\end{itemize}

\begin{tcolorbox}[breakable,colback=teal!4,colframe=teal!35!black,title={Before you move on}]
What does circa mean? (\textquotedblleft{}About, approximately\textquotedblright{}.) Say it once more.
\end{tcolorbox}
\section[soltanto]{soltanto — only}

\label{lesson:IT-C244-soltanto}



\begin{quote}
Before the new one: say the Italian for almost, then the Italian for about.
\end{quote}

\subsection*{You'll want to know: soltanto}
\textbf{soltanto} — \textquotedblleft{}only\textquotedblright{}.

\begin{cousinweb}
\textbf{soltanto} is \emph{solo}, \textquotedblleft{}alone, only\textquotedblright{}, plus \emph{tanto}, \textquotedblleft{}so much\textquotedblright{}: only so much. \emph{solo} $\leftarrow$ Latin \emph{solus}, \textquotedblleft{}alone\textquotedblright{}; \emph{tanto} $\leftarrow$ Latin \emph{tantus}. (EtymDB 2.1)
\end{cousinweb}

\begin{grammarlens}[title={What you've built}]
One more word for this chapter.
\end{grammarlens}

\subsection*{Your turn}
\begin{itemize}
\raggedright
  \item \emph{Say it:} \emph{soltanto}
  \item \emph{Say it:} \emph{soltanto}, once more
  \item \emph{Say it:} say \emph{soltanto}
  \item \emph{Recall:} say the Italian for almost, then the Italian for about, then say \emph{soltanto} again
\end{itemize}

\begin{tcolorbox}[breakable,colback=teal!4,colframe=teal!35!black,title={Before you move on}]
What does soltanto mean? (\textquotedblleft{}Only\textquotedblright{}.) Say it once more.
\end{tcolorbox}
\section[insieme]{insieme — together}

\label{lesson:IT-C244-insieme}



\begin{quote}
Before the new one: say the Italian for about, then the Italian for only.
\end{quote}

\subsection*{You'll want to know: insieme}
\textbf{insieme} — \textquotedblleft{}together\textquotedblright{}.

\begin{cousinweb}
\textbf{insieme} $\leftarrow$ Latin \emph{insimul}, \textquotedblleft{}at the same time\textquotedblright{} (Lewis \& Short) — inherited; EtymDB 2.1 has no row for \emph{insieme}.
\end{cousinweb}

\begin{grammarlens}[title={What you've built}]
One more word for this chapter.
\end{grammarlens}

\subsection*{Your turn}
\begin{itemize}
\raggedright
  \item \emph{Say it:} \emph{insieme}
  \item \emph{Say it:} \emph{insieme}, once more
  \item \emph{Say it:} say \emph{insieme}
  \item \emph{Recall:} say the Italian for about, then the Italian for only, then say \emph{insieme} again
\end{itemize}

\begin{tcolorbox}[breakable,colback=teal!4,colframe=teal!35!black,title={Before you move on}]
What does insieme mean? (\textquotedblleft{}Together\textquotedblright{}.) Say it once more.
\end{tcolorbox}
